% =============================================================================
% Project : THE_PHILOSOPHY_OF_VIABILITY
% Book    : Book I — Ontology of Viability
% Book ID : VS.B01
%
% File    : 4_ASSETS/diagrams/
%           VS.B01.P01.C02.S09.SS02.P004.DG01.forms_and_domains_of_invariance.tex
%
% Asset   : DG01 — Forms and Domains of Invariance
% Location: Chapter 02 / Section 09 / Subsection 02 / Paragraph 004
%
% Purpose :
%   To visualize the analytical independence of forms of
%   invariance and domains of invariance.
%
% Conceptual Caution :
%   The figure represents possible analytical cross-classification.
%   It does not assert that every combination is instantiated
%   or ontologically established.
%
% Requires (preamble):
%   \usepackage{tikz}
%   \usetikzlibrary{calc, positioning, fit}
% =============================================================================

\begin{figure}[htbp]
    \centering

    \begin{tikzpicture}[
        font=\small,
        node distance=0.4cm,
        % --- base box ---
        box/.style={
            draw, rounded corners=3pt, align=center,
            inner sep=5pt, line width=0.6pt
        },
        % --- node roles (names avoid clashes with built-in TikZ keys) ---
        formbox/.style={box, text width=3.4cm, minimum height=2.3cm,
                        fill=blue!8, draw=blue!55!black},
        domainbox/.style={box, text width=2.4cm, minimum height=1.4cm,
                          fill=gray!10, draw=black!60},
        % --- frames, titles, labels ---
        vsgroup/.style={draw=black!45, dashed, rounded corners=5pt,
                        inner sep=8pt, line width=0.5pt},
        axistitle/.style={font=\small\bfseries, align=center},
        grouplabel/.style={font=\footnotesize\itshape, text=black!75,
                           align=center},
        % --- undirected link (no arrowheads: no priority between axes) ---
        vslink/.style={draw=black!50, dashed, line width=0.6pt}
    ]

    % ---------------------------------------------------------------
    % Form axis (left)
    % ---------------------------------------------------------------
    \node[formbox] (question) {
        What determination is claimed to remain preserved?
    };
    \node[formbox, below=of question] (examples) {
        Structural, relational, quantitative, qualitative,
        compositional, and other analytical forms
    };

    % ---------------------------------------------------------------
    % Domain axis (right): two columns of three, top-aligned with Form
    % ---------------------------------------------------------------
    \node[domainbox, anchor=north west] (d1) at
        ($(question.north east)+(1.8,0)$) {Formal Domain};
    \node[domainbox, below=of d1] (d2) {First Determination};
    \node[domainbox, below=of d2] (d3) {Structural Determination};

    \node[domainbox, right=of d1] (d4) {Relational Manifestation};
    \node[domainbox, below=of d4] (d5) {Material Manifestation};
    \node[domainbox, below=of d5] (d6) {Intermediate Hubs};

    % ---------------------------------------------------------------
    % Frames, titles, labels
    % ---------------------------------------------------------------
    \node[vsgroup, fit=(question)(examples)] (groupF) {};
    \node[vsgroup, fit=(d1)(d2)(d3)(d4)(d5)(d6)] (groupD) {};

    \node[axistitle, above=0.2cm of groupF] {FORM OF INVARIANCE};
    \node[axistitle, above=0.2cm of groupD] {DOMAIN OF EXAMINATION};

    \node[grouplabel, below=0.15cm of groupD] {
        Distinct domains of examination
    };

    % ---------------------------------------------------------------
    % Independence marker (symmetric, no arrowheads)
    % ---------------------------------------------------------------
    \node[circle, draw=black!60, fill=white, inner sep=2pt, font=\large]
        (times) at ($(groupF.east)!0.5!(groupD.west)$) {$\times$};
    \draw[vslink] (groupF.east) -- (times);
    \draw[vslink] (times) -- (groupD.west);

    % ---------------------------------------------------------------
    % Note (centred under both axes)
    % ---------------------------------------------------------------
    \coordinate (mid) at ($(groupF.south)!0.5!(groupD.south)$);
    \node[below=0.9cm of mid, text width=10.5cm, align=center,
          font=\footnotesize\itshape] {
        The form identifies the respect of preservation;
        the domain identifies the field in which that
        preservation is examined. Neither dimension
        determines the other.
    };

    \end{tikzpicture}

    \caption{Forms and domains constitute distinct analytical
    dimensions. The diagram presents their independence without
    asserting that every possible form-domain combination has
    been established.}
    \label{fig:forms-and-domains-of-invariance}
\end{figure}

% =============================================================================
% END OF DIAGRAM DG01
% =============================================================================